\documentclass[12pt]{article}
\usepackage[a4paper, margin=1in]{geometry}
\usepackage{graphicx}
\usepackage{booktabs}
\usepackage{longtable}
\usepackage{amsmath}
\usepackage{amssymb}
\usepackage{float}
\usepackage{setspace}
\usepackage{indentfirst}
\usepackage{hyperref}
\hypersetup{
    colorlinks=true,
    linkcolor=black,
    filecolor=black,
    urlcolor=black,
    citecolor=black,
}
\usepackage{tabularx}
\usepackage{array}
\usepackage{caption}
\usepackage{subcaption}
\usepackage{microtype}
\usepackage{natbib}

\title{\textbf{Modeling and Forecasting China's Electric Vehicle Sales\\
Using a Seasonally-Optimized SARIMA Model}}

\author{Ngo Hoang Phuc \\[4pt]
National Economics University \\
\texttt{11230582@st.neu.edu.vn}}

\date{May 2026}

\begin{document}
\onehalfspacing
\maketitle

\begin{abstract}
\noindent
This paper develops a Seasonal Autoregressive Integrated Moving Average (SARIMA) framework
to model and forecast monthly electric vehicle (EV) sales in China from January 2018 to
April 2024 (76 observations). The dataset is constructed from the China Automobile Sales
Data (Kaggle) by filtering EV records and aggregating to monthly totals. After thorough
preprocessing, exploratory analysis, and decomposition, we confirm that the series
exhibits a strong long-term upward trend, pronounced 12-month seasonality, and
non-constant variance driven by rapid market growth. Stationarity is verified jointly
using the Augmented Dickey-Fuller (ADF) and KPSS tests, which indicate that
first-order non-seasonal differencing ($d=1$) combined with one seasonal difference
($D=1, s=12$) renders the series fully stationary.

Model selection proceeds in two stages: an initial stepwise auto-ARIMA search guided
by the Bayesian Information Criterion (BIC), followed by a nested rolling
cross-validation procedure that selects the final specification
SARIMA$(1,1,0)\times(1,1,0)_{12}$ based on minimum out-of-sample RMSE across multiple
time folds. On a held-out 12-month test window (May 2023 to April 2024), the model
achieves RMSE~$=108{,}191$ units and MAPE~$=17.42\%$, with Ljung-Box residual test
$p=0.336$, confirming white-noise residuals. SARIMA outperforms a non-seasonal ARIMA
and a simple LSTM, and is competitive with Holt-Winters and Prophet. Robustness checks
confirm that modeling the original series without log transformation yields lower errors
than a multiplicative log-space specification.

The 12-month forward forecast (May 2024--April 2025) indicates continued growth, with
sales peaking around December 2024 at approximately 810,000 units, followed by an
early-year dip to roughly 446,000 units in February 2025 due to the Lunar New Year
holiday effect. These projections provide actionable guidance for production planning,
inventory management, supply chain coordination, and public charging infrastructure
deployment.
\end{abstract}

\vspace{0.3cm}
\noindent \textbf{Keywords:} Time Series Analysis, SARIMA Model, Electric Vehicles,
Demand Forecasting, Seasonality, China NEV Market.\\
\noindent \textbf{JEL Classification:} C22, Q47, L62.

\newpage
\tableofcontents
\newpage

%% ─────────────────────────────────────────────────────────────────────────────
\section{Introduction}
\label{sec:intro}
%% ─────────────────────────────────────────────────────────────────────────────

The electric vehicle market has become one of the most strategically important segments
of the global automotive industry. Driven by a convergence of environmental policy,
technological advancement, and shifting consumer preferences, global EV sales have grown
exponentially over the past decade. China occupies the commanding position in this
expansion: it is simultaneously the world's largest producer and largest consumer of
new energy vehicles (NEVs), accounting for more than half of global EV sales annually.
As of December 2023, monthly EV sales in China reached approximately 756,000 units---a
figure that was unimaginable only six years earlier, when monthly sales stood below
25,000 units in early 2018.

This extraordinary growth trajectory creates both an opportunity and a challenge for
forecasters. The market is not merely expanding---it is doing so with pronounced
seasonality, episodic policy-driven shocks, and structural breaks that complicate
standard time series models. Understanding and anticipating these dynamics is critical
for a wide range of stakeholders. Automotive manufacturers must plan production
schedules and allocate battery procurement well in advance of demand peaks.
Distributors and retailers need to build inventory buffers aligned with predictable
seasonal cycles. Policymakers responsible for charging infrastructure deployment require
reliable forward-looking demand estimates to site and scale public charging networks.
Financial investors and analysts need forecasts to calibrate valuation models and
assess market saturation risk.

Despite the practical importance of accurate EV sales forecasting in China, the
existing academic literature is still relatively thin on this specific topic,
particularly with respect to transparent, reproducible methodologies grounded in
rigorous statistical testing. Several recent studies have applied SARIMA-family models
to EV sales in China and the United States with promising results, but they differ in
their data coverage, model selection strategies, and validation approaches
\citep{chen2025, cetin2024}. This paper seeks to fill part of that gap by providing a
complete, end-to-end SARIMA forecasting workflow applied to a long-horizon China EV
panel.

The specific contributions of this paper are threefold. First, it constructs a clean,
monthly China EV sales series spanning January 2018 to April 2024 by aggregating
vehicle-level transaction records, documenting all preprocessing steps for
reproducibility. Second, it implements a disciplined two-stage model selection
procedure that combines auto-ARIMA with nested rolling cross-validation
(rolling CV)---a more robust alternative to single-split evaluation. Third, it provides
a systematic benchmark comparison against Holt-Winters, Prophet, non-seasonal ARIMA,
and a simple LSTM, and discusses the economic interpretability of the chosen model's
forecasts in the context of China's NEV market dynamics.

The remainder of the paper is organized as follows. Section~\ref{sec:lit} reviews
relevant prior literature. Section~\ref{sec:data} describes the data and preprocessing
steps, including exploratory analysis and decomposition. Section~\ref{sec:methods}
presents the econometric framework, including the SARIMA model specification, stationarity
testing, and the nested CV selection strategy.
Section~\ref{sec:results} reports empirical results covering stationarity tests, model
selection, residual diagnostics, forecast accuracy, and a comparative assessment against
baseline models. Section~\ref{sec:discussion} discusses the economic implications,
robustness, and limitations of the results. Section~\ref{sec:conclusion} concludes.

%% ─────────────────────────────────────────────────────────────────────────────
\section{Literature Review}
\label{sec:lit}
%% ─────────────────────────────────────────────────────────────────────────────

Forecasting automobile sales is a long-standing topic in applied econometrics, drawing
on methods ranging from structural demand models to reduced-form time series techniques.
Within the electric vehicle literature, forecasting research has grown substantially
since around 2020, motivated by the dramatic acceleration of market adoption and its
policy implications.

\subsection{Seasonal Time Series Modeling for Demand Forecasting}

The ARIMA family of models, first systematized by Box and Jenkins, remains a
foundational approach for univariate time series with trend and autocorrelation
structure \citep{cetin2024}. For series exhibiting periodic seasonal cycles, the
Seasonal ARIMA (SARIMA) extension adds seasonal AR, differencing, and MA terms,
enabling the model to capture recurring within-year patterns simultaneously with
non-seasonal dynamics. This property makes SARIMA particularly attractive for monthly
automotive sales data, where strong intra-year variation driven by holidays, promotions,
and policy cycles is well-documented.

The comparative evidence on SARIMA's performance relative to more flexible machine
learning alternatives is nuanced. \citet{cetin2024} summarize a cross-domain literature
including studies on electricity load, rainfall, unemployment, retail demand, and
passenger traffic. Their review finds that SARIMA and hybrid SARIMA-GARCH models
consistently achieve high accuracy for data governed by seasonal patterns, while deep
learning approaches such as LSTM and ANN tend to outperform only in settings with
abundant high-frequency data. For monthly series with fewer than 100 observations---the
regime applicable to this paper---statistical parsimony favors SARIMA.

\subsection{EV Sales Forecasting in China and the United States}

\citet{chen2025} apply a SARIMA framework specifically to monthly new energy vehicle
(NEV) sales in China from January 2022 to April 2024. Their study documents strong
seasonality, with fourth-quarter peaks attributable to year-end promotions and
policy-induced buying acceleration. They demonstrate that SARIMA outperforms
Holt-Winters for this particular dataset and argue that the model's transparency is
valuable for downstream applications such as charging infrastructure planning. One
important finding is that policy events---particularly subsidy changes---generate abrupt
demand shifts that univariate models struggle to anticipate, motivating extensions to
SARIMAX with exogenous policy variables.

\citet{cetin2024} forecast monthly US light-duty EV sales for 2024--2025 using an
optimized SARIMA model selected via grid search over AIC, BIC, and HQIC criteria.
Their preferred specification, SARIMA$(1,0,2)\times(1,0,1)_{12}$, predicts a
consistent upward trend in US EV adoption, with seasonal peaks in summer months
(June--July) reflecting stronger warm-weather consumer activity. These results
stand in contrast to China's pattern, where peaks typically concentrate in Q4
under year-end promotional and subsidy dynamics---a market-structure difference that
motivates country-specific modeling.

\subsection{Alternative Forecasting Methods}

Beyond SARIMA, several alternative approaches have been applied in related contexts.
The Holt-Winters exponential smoothing (ETS) model adds level, trend, and seasonal
smoothing components and is computationally simpler but lacks the ACF/PACF-based
structural interpretation of SARIMA \citep{cetin2024}. Prophet, developed at Facebook,
uses an additive regression model with trend changepoints and Fourier-series seasonality;
it can be powerful for series with multiple seasonal layers or irregular trends but has
been shown to underperform SARIMA on short monthly series without complex multi-frequency
patterns \citep{cetin2024}. LSTM neural networks offer flexibility for high-frequency
data but are data-hungry, and several studies confirm they cannot reliably outperform
classical methods on monthly panels with fewer than 100 observations.

\subsection{Gap Addressed by This Paper}

The literature has three notable gaps that this paper addresses. First, existing
China-focused EV forecasting studies use either short data windows (2022--2024) or
aggregate data that smooth over micro-level heterogeneity; this paper uses a longer
series from 2018 that captures multiple market phases. Second, most studies rely on
a single hold-out split for evaluation, which may yield unreliable performance
estimates when the test period coincides with a market shock; this paper instead uses
nested rolling CV. Third, explicit model robustness checks---such as
comparing additive versus multiplicative specifications---are rarely conducted in the EV
forecasting literature, yet are important for choosing the correct model family.
Table~\ref{tab:lit} summarizes selected related studies and highlights the positioning
of this work.

\begin{table}[H]
\centering
\small
\caption{Selected literature on EV and seasonal demand forecasting}
\label{tab:lit}
\begin{tabular}{p{3.0cm} p{2.5cm} p{3.0cm} p{5.2cm}}
\toprule
Study & Method & Data scope & Key finding / contribution \\
\midrule
\citet{chen2025} & SARIMA & China NEV, 2022--2024 & Strong Q4 seasonality; policy-induced structural breaks \\[4pt]
\citet{cetin2024} & SARIMA (grid search) & US light-duty EV, 2021--2023 & Summer peaks; AIC/BIC selection; upward trend confirmed \\[4pt]
\citeauthor{cetin2024}\ (\citeyear{cetin2024}), survey & SARIMA vs.\ LSTM/ANN & Multiple domains & SARIMA competitive on short seasonal series \\[4pt]
This paper & SARIMA (nested CV) & China EV, 2018--2024 & Longer panel; robust nested CV selection; baseline comparisons \\
\bottomrule
\end{tabular}
\end{table}

%% ─────────────────────────────────────────────────────────────────────────────
\section{Data and Exploratory Analysis}
\label{sec:data}
%% ─────────────────────────────────────────────────────────────────────────────

\subsection{Data Source}

The primary data source is the China Automobile Monthly Sales Data (2018--2024.4)
made available on Kaggle by Felixzhao \citep{felixzhao2024}. The raw dataset contains
38,806 vehicle-model-month observations covering a broad range of passenger vehicle
models sold in China, with fields including \texttt{year\_month}, \texttt{units\_sold},
\texttt{brand}, \texttt{model}, and \texttt{is\_ev}, where the last indicator
distinguishes electric vehicles from internal combustion engine (ICE) models. The data
are compiled from official sales registration records reported by the China Passenger Car
Association (CPCA) and partner market research firms, ensuring coverage of national
market transactions.

\subsection{Data Preprocessing}

The preprocessing pipeline proceeds in five steps. First, \textit{deduplication}: a
review identified 33 fully duplicated rows, which were removed to prevent double-counting
in the monthly aggregation. Second, \textit{type conversion}: the \texttt{units\_sold}
field was stored as a string with comma separators; these were stripped and the field
converted to integer. No negative or non-numeric values remained after conversion.
Third, \textit{EV filtering}: only rows where \texttt{is\_ev == ``EV''} were retained.
Fourth, \textit{monthly aggregation}: EV unit counts were summed within each calendar
month, yielding a monthly total-sales series indexed by month-start dates. Fifth,
\textit{frequency alignment}: the index was set to a regular monthly start (\texttt{MS})
frequency, and any missing months were filled with zero (no missing months were found
in practice). The resulting processed series is exported to
\texttt{ev\_sales\_monthly.csv} and covers January 2018 to April 2024, for a total of
$T = 76$ observations.

Formally, the monthly EV sales series is defined as:
\begin{equation}
\text{EV\_Sales}_{t} = \sum_{i \in \mathcal{M}(t)} \text{units\_sold}_{i},
\end{equation}
where $\mathcal{M}(t)$ is the set of all EV model-records belonging to calendar month $t$.

\subsection{Descriptive Statistics}

Table~\ref{tab:desc} reports summary statistics for the processed series.

\begin{table}[H]
\centering
\caption{Summary statistics for monthly China EV sales (2018-01 to 2024-04)}
\label{tab:desc}
\begin{tabular}{lr}
\toprule
Statistic & Value \\
\midrule
Observations & 76 \\
Mean & 255,718 \\
Standard deviation & 214,734 \\
Minimum & 8,988 (Feb 2020) \\
25th percentile & 68,143 \\
Median & 168,863 \\
75th percentile & 425,222 \\
Maximum & 755,943 (Dec 2023) \\
\bottomrule
\end{tabular}
\end{table}

The statistics underscore the series' exceptional dynamism. The coefficient of
variation exceeds 80\%, reflecting not merely noise but a fundamental structural
expansion of the market from fewer than 25,000 units per month in early 2018 to
almost 756,000 units in December 2023---a 30-fold increase in six years.

\subsection{Visual Inspection and Trend Analysis}

The raw time series (Figure~\ref{fig:series}) reveals three distinct growth phases.
The first phase (2018--2020) exhibits modest, low-level growth punctuated by a
severe trough in February 2020 (8,988 units), which corresponds to the national
COVID-19 lockdowns that halted automotive production and retail activity. The second
phase (2021--2022) represents explosive market acceleration: monthly sales surged
from roughly 97,000 units in February 2021 to over 651,000 units in December 2022.
This acceleration was powered by aggressive government purchase subsidies, rapid
expansion of the charging network, and the entry of new domestic OEM brands
(notably BYD's dramatic commercial rise). The third phase (2023--2024) is
characterized by continued growth at a more moderate, stabilizing rate following the
complete phase-out of national EV purchase subsidies on December 31, 2022.

\begin{figure}[H]
\centering
\includegraphics[width=0.92\textwidth]{figures/02_ev_sales_monthly.png}
\caption{Monthly EV sales in China, January 2018 to April 2024. The sharp trough
in February 2020 reflects COVID-19 lockdowns; the sharp dip in January 2023
reflects the ``post-subsidy hangover'' following the expiry of national EV
subsidies on 31 December 2022.}
\label{fig:series}
\end{figure}

\subsection{Seasonal Decomposition}

To separate the structural components of the series, we apply additive seasonal
decomposition with a 12-month period. The decomposition reveals:

\begin{enumerate}
\item \textbf{Trend:} A smooth, accelerating upward trend from 2018 through 2022,
followed by a leveling off in 2023--2024 as the market matures and the subsidy
environment normalizes.

\item \textbf{Seasonality:} A stable 12-month seasonal pattern with consistent
within-year amplitude, ranging from approximately $-$75,000 units (February trough)
to $+$75,000 units (November--December peak). The January--February dip is attributable
to the Lunar New Year holiday effect, which suppresses consumer activity and disrupts
manufacturing. The late-year surge reflects year-end promotions, fleet purchases, and
historically, the rush to capitalize on expiring subsidies.

\item \textbf{Residuals:} Residuals are well-behaved from 2018 through most of 2022.
A large negative outlier emerges in January 2023 (approximately $-$175,000 units
relative to the trend-seasonal component), corresponding to the post-subsidy
demand collapse. A 3-sigma outlier test confirms this as the single statistically
anomalous observation in the sample.
\end{enumerate}

\begin{figure}[H]
\centering
\includegraphics[width=0.92\textwidth]{figures/04_seasonal_decompose.png}
\caption{Additive seasonal decomposition of monthly EV sales: observed series (top),
trend component (second), seasonal component (third), and residuals (bottom).
The seasonal component shows a stable 12-month cycle; the residuals reveal a
large negative shock in January 2023.}
\label{fig:decomp}
\end{figure}

We also compare additive and multiplicative decomposition models. The multiplicative
residuals have a mean closer to 1.0 (ideal) on the training data, suggesting
multiplicative seasonal structure. However, as demonstrated in Section~\ref{sec:robust},
modeling the series directly in levels under an additive SARIMA specification
outperforms a log-space (multiplicative) SARIMA on the out-of-sample test window.
This apparent contradiction is explained by the structural break in 2023: the
multiplicative specification amplifies the early-period growth trajectory and
over-extrapolates after the market shift, while the additive model---being more
conservative---adapts more gracefully to the deceleration phase.

The January 2023 outlier is retained without adjustment. It reflects a genuine
policy-driven demand shock rather than a measurement error, and removing it would
discard economically meaningful information about structural market dynamics.

%% ─────────────────────────────────────────────────────────────────────────────
\section{Methodology}
\label{sec:methods}
%% ─────────────────────────────────────────────────────────────────────────────

\subsection{The SARIMA Model}

The Seasonal Autoregressive Integrated Moving Average (SARIMA) model generalizes
the Box-Jenkins ARIMA framework to accommodate periodic seasonality. The general
SARIMA$(p,d,q)\times(P,D,Q)_{s}$ model is:
\begin{equation}
\Phi_{P}(B^{s})\,\phi_{p}(B)\,\nabla_{s}^{D}\nabla^{d}\,x_{t}
= \Theta_{Q}(B^{s})\,\theta_{q}(B)\,w_{t},
\label{eq:sarima}
\end{equation}
where $B$ is the backshift operator ($Bx_t = x_{t-1}$), $\nabla^d = (1-B)^d$ is
the non-seasonal difference operator, $\nabla_s^D = (1-B^s)^D$ is the seasonal
difference operator, $\phi_p(B) = 1 - \phi_1 B - \cdots - \phi_p B^p$ and
$\theta_q(B) = 1 + \theta_1 B + \cdots + \theta_q B^q$ are the non-seasonal AR and
MA polynomials, $\Phi_P(B^s)$ and $\Theta_Q(B^s)$ are their seasonal counterparts at
lag $s$, and $w_t \sim \text{WN}(0, \sigma^2)$ is a white noise innovation process
\citep{cetin2024}.

The non-seasonal parameters $(p, d, q)$ govern short-term autocorrelation structure
after trend removal, while the seasonal parameters $(P, D, Q)$ govern year-over-year
dynamics at period $s = 12$ months. For China EV sales, the 12-month period is
economically motivated by the annual promotional cycle, the Lunar New Year effect,
and policy announcement calendars.

\subsection{Stationarity Testing}

Consistent parameter estimation in ARIMA-family models requires stationarity of the
transformed series. We conduct a dual testing strategy using:

\begin{itemize}
\item \textbf{Augmented Dickey-Fuller (ADF) test:} Null hypothesis is a unit root
(non-stationarity). Rejection ($p < 0.05$) supports stationarity.
\item \textbf{KPSS test:} Null hypothesis is stationarity. Failure to reject
($p > 0.05$) supports stationarity.
\end{itemize}

Using both tests in conjunction reduces the risk of a false conclusion. An ideal
outcome is simultaneous ADF rejection and KPSS non-rejection \citep{cetin2024}.

We evaluate differencing combinations $(d, D) \in \{0,1,2\} \times \{0,1\}$ and
select the minimum transformation that achieves confirmed stationarity. The seasonal
period is fixed at $s = 12$ throughout.

\subsection{Model Selection: Auto-ARIMA and Nested Rolling Cross-Validation}

Model selection follows a two-stage procedure designed to balance fit quality with
generalization:

\textbf{Stage 1 --- Auto-ARIMA (BIC-guided):} An initial stepwise search over
$(p, q) \in \{0,1,2\}^2$ and $(P,Q) \in \{0,1\}^2$ is conducted using the
\texttt{pmdarima} library with the differencing orders fixed at $(d,D)=(1,1)$ and
seasonal period $s=12$. The Bayesian Information Criterion (BIC) is used as the
selection criterion because it penalizes model complexity more heavily than AIC,
reducing the risk of over-parameterization in a sample of only 64 training
observations.

\textbf{Stage 2 --- Nested Rolling Cross-Validation:} The auto-ARIMA output provides
an initial candidate set. To obtain a robust out-of-sample evaluation, we apply nested
rolling cross-validation (nested CV): for each rolling fold (expanding window with
step size 1 and horizon 1), a fresh auto-ARIMA is fitted on the fold's training data,
and the one-step forecast is compared to the held-out observation. The configuration
that achieves the lowest mean RMSE across all folds is selected as the final model.

This nested CV approach is motivated by the literature on temporal cross-validation,
which shows that a single hold-out split can yield misleading performance estimates
when the data exhibit structural breaks \citep{cetin2024}. In the presence of the
January 2023 policy shock, the nested CV is particularly valuable because it evaluates
the model's robustness across different historical windows.

\subsection{Train-Test Split and Final Evaluation}

The 76-observation series is split into a training set of 64 months (January 2018
to April 2023) and a test set of 12 months (May 2023 to April 2024). This 64/12 split
places the test window in the post-subsidy market regime, providing a demanding
out-of-sample evaluation that tests whether the model can generalize to a structural
environment different from the training period.

Forecast accuracy on the test window is measured by:
\begin{equation}
\text{RMSE} = \sqrt{\frac{1}{n}\sum_{t=1}^{n}(y_t - \hat{y}_t)^2},
\qquad
\text{MAPE} = \frac{1}{n}\sum_{t=1}^{n}\left|\frac{y_t - \hat{y}_t}{y_t}\right|
\times 100\%.
\end{equation}
Residual independence is assessed using the Ljung-Box $Q$-statistic at lag $h = 12$:
\begin{equation}
Q(12) = n(n+2)\sum_{j=1}^{12}\frac{\hat{\rho}_j^2}{n-j},
\end{equation}
where $\hat{\rho}_j$ is the sample autocorrelation of residuals at lag $j$. Under
the null hypothesis of white-noise residuals, $Q(12) \sim \chi^2(12)$.

\subsection{Baseline Models}

To contextualize the SARIMA performance, four baseline models are estimated on the
same training set and evaluated on the same test split:

\begin{itemize}
\item \textbf{Non-seasonal ARIMA:} auto-ARIMA with \texttt{seasonal=False},
capturing only trend dynamics.
\item \textbf{Holt-Winters (ETS):} Additive trend and additive seasonal smoothing
with 12-month seasonal period.
\item \textbf{Prophet:} Facebook's additive regression model with Fourier-series
annual seasonality.
\item \textbf{LSTM:} A simple one-layer LSTM with 32 hidden units, MinMax
normalization, and lookback windows of 6, 12, and 24 months; the best configuration
is reported.
\end{itemize}

\subsection{Robustness Check: Additive vs.\ Multiplicative (Log-Space)}

To verify that the additive specification is preferable to a multiplicative
alternative, we replicate the full nested-CV pipeline on $\log(1 + y_t)$ and
back-transform forecasts to the original scale for comparison. A lower test RMSE
under the no-log specification confirms the additive choice.

\subsection{Final Forecast}

After validation, the model is re-estimated on the full 76-observation series and
used to generate a 12-month forward forecast (May 2024 to April 2025) with 95\%
confidence intervals. The confidence interval width reflects the variance of the
innovation process and widens with the forecast horizon in the standard manner.

%% ─────────────────────────────────────────────────────────────────────────────
\section{Results}
\label{sec:results}
%% ─────────────────────────────────────────────────────────────────────────────

\subsection{ACF and PACF Analysis}

Figure~\ref{fig:acf_levels} shows the ACF and PACF of the raw EV sales series.
The ACF decays very slowly and almost linearly across 24 lags, with all coefficients
far outside the 95\% confidence band. The PACF shows a single dominant spike at
lag 1 (close to 1.0), followed by rapid decay. This pattern is the canonical
signature of a strongly non-stationary series with a persistent trend component
\citep{cetin2024}.

\begin{figure}[H]
\centering
\includegraphics[width=0.92\textwidth]{figures/05_acf_pacf_raw.png}
\caption{ACF (left) and PACF (right) of monthly EV sales in levels. The slowly
decaying ACF confirms non-stationarity; the dominant PACF spike at lag 1 is
consistent with a near-unit-root process.}
\label{fig:acf_levels}
\end{figure}

After applying combined non-seasonal and seasonal differencing $(d=1, D=1, s=12)$,
the transformed series shows dramatically improved behavior (Figure~\ref{fig:acf_diff}).
The ACF exhibits a sharp negative spike at lag 1 and then drops within the confidence
band, while the PACF shows a similar pattern. This structure is consistent with a
low-order ARMA process, confirming that the chosen differencing is sufficient and
appropriate.

\begin{figure}[H]
\centering
\includegraphics[width=0.92\textwidth]{figures/06_acf_pacf_diff.png}
\caption{ACF (left) and PACF (right) after differencing $(d=1, D=1, s=12)$.
The series is now effectively stationary, with residual autocorrelations
largely within the 95\% confidence band.}
\label{fig:acf_diff}
\end{figure}

\subsection{Stationarity Tests}

Table~\ref{tab:stationarity} reports the ADF and KPSS test results for the raw and
differenced series.

\begin{table}[H]
\centering
\caption{Stationarity tests: ADF and KPSS results}
\label{tab:stationarity}
\begin{tabular}{llcccc}
\toprule
Series & $(d, D)$ & ADF stat. & ADF $p$-value
& KPSS stat. & KPSS $p$-value \\
\midrule
Raw data & $(0, 0)$ & $-1.012$ & $0.9579$ & $1.2073$ & $0.010$ \\
Differenced & $(1, 1)$ & $-12.375$ & $\approx 0$ & $0.1284$ & $0.100$ \\
\bottomrule
\end{tabular}
\end{table}

The raw series fails both stationarity criteria: the ADF $p$-value of 0.9579 cannot
reject the unit root null, and the KPSS statistic of 1.2073 far exceeds the 5\%
critical value of 0.463, strongly rejecting the stationarity null. After applying
$(d=1, D=1)$ differencing, both tests confirm stationarity with high confidence:
the ADF statistic of $-12.375$ is well below the 1\% critical value ($-3.449$),
and the KPSS statistic of 0.1284 falls safely below the 10\% critical value (0.347),
so the stationarity null is not rejected. This dual confirmation---ADF rejection and
KPSS non-rejection---is the gold standard for stationarity verification
\citep{cetin2024}. We therefore fix $d=1$ and $D=1$ for all subsequent SARIMA
estimation.

\subsection{Model Selection and the Final Specification}

\textbf{Stage 1 --- Auto-ARIMA (BIC):} The stepwise auto-ARIMA search over the
constrained parameter grid ($p, q \le 2$; $P, Q \le 1$; $d=1$; $D=1$; $s=12$)
identifies SARIMA$(2,1,2)\times(1,1,1)_{12}$ as the BIC-optimal specification on
the 64-month training set, with BIC = 910.67 and AIC = 899.59. Several competing
specifications are close in information criteria; Table~\ref{tab:aic_bic} reports
the top candidates.

\begin{table}[H]
\centering
\caption{Top SARIMA candidate models by AIC/BIC (training set)}
\label{tab:aic_bic}
\begin{tabular}{lcc}
\toprule
Model & AIC & BIC \\
\midrule
$(2,1,2)\times(1,1,1)_{12}$ & 899.59 & 910.67 \\
$(3,1,1)\times(1,1,1)_{12}$ & 900.65 & 911.74 \\
$(1,1,2)\times(0,1,1)_{12}$ & 911.43 & 919.35 \\
$(2,1,1)\times(1,1,0)_{12}$ & 935.41 & 943.47 \\
$(1,1,1)\times(0,1,1)_{12}$ & 936.61 & 943.06 \\
\bottomrule
\end{tabular}
\end{table}

\textbf{Stage 2 --- Nested Rolling CV:} The nested rolling cross-validation
procedure, run over expanding folds with an initial window of at least 24
observations, selects \textbf{SARIMA$(1,1,0)\times(1,1,0)_{12}$} as the model with
the lowest mean RMSE across all folds ($\text{Mean RMSE}_{\text{CV}} \approx
49{,}561$; $\text{Mean MAPE}_{\text{CV}} \approx 17.08\%$). This simpler
specification appeared as the optimal choice in 6 out of the nested folds, indicating
strong stability and adaptability as the training window expands over time.

The nested-CV result differs from the BIC-optimal auto-ARIMA model because BIC
evaluates in-sample fit, while nested CV evaluates rolling out-of-sample
performance. The fact that a simpler model wins on the out-of-sample criterion is
consistent with the bias-variance tradeoff: the more heavily parameterized
$(2,1,2)\times(1,1,1)$ specification overfits the training data, whereas the
parsimonious $(1,1,0)\times(1,1,0)$ generalizes better to unseen horizons.

The two estimated coefficients in the final model are both statistically significant:
the non-seasonal AR term ($\text{ar.L1} = -0.322$, $p = 0.038$) captures a
mean-reversion tendency at the one-month lag, and the seasonal AR term
($\text{ar.S.L12} = -0.564$) reflects a negative dependency on same-month sales
one year earlier, consistent with the market's year-over-year deceleration pattern
after 2022.

\subsection{Residual Diagnostics}

Figure~\ref{fig:resid} presents four diagnostic panels for the selected model.

\begin{figure}[H]
\centering
\includegraphics[width=0.92\textwidth]{figures/07_residual_diagnostics.png}
\caption{Residual diagnostics for SARIMA$(1,1,0)\times(1,1,0)_{12}$:
residuals over time (top left), histogram (top right), residual ACF (bottom left),
and Q-Q plot (bottom right). Most autocorrelations lie within the 95\% confidence
band, and the Ljung-Box $p$-value at lag 12 is 0.3355, confirming white-noise
residuals.}
\label{fig:resid}
\end{figure}

The residuals fluctuate around zero with no obvious trend from 2018 through mid-2022.
A cluster of large negative residuals appears in late 2022 to early 2023, reflecting
the post-subsidy demand collapse that the univariate model cannot capture with
lagged information alone. The histogram is approximately bell-shaped and symmetric
around zero, with mild excess kurtosis in the left tail from the 2022--2023 shock.
The residual ACF shows no significant spikes beyond lag 0, and the Q-Q plot aligns
well with the theoretical normal line in the middle quantile range, with expected
heavy-tail deviations at the extremes.

The Ljung-Box $Q$-statistic at lag 12 yields $Q(12) = 14.91$ with a $p$-value of
$0.3355 > 0.05$, confirming that the model residuals are not serially correlated.
This satisfies the core adequacy condition for SARIMA and validates the model
specification.

\subsection{Test Set Forecast Performance}

Table~\ref{tab:test_metrics} reports forecast accuracy on the 12-month test window.

\begin{table}[H]
\centering
\caption{Test set performance: SARIMA$(1,1,0)\times(1,1,0)_{12}$ (May 2023 -- April 2024)}
\label{tab:test_metrics}
\begin{tabular}{lcc}
\toprule
Split & RMSE (units) & MAPE (\%) \\
\midrule
Training set & 55,014 & 28.56 \\
Test set & 108,191 & 17.42 \\
Nested CV (mean across folds) & 49,561 & 17.08 \\
\bottomrule
\end{tabular}
\end{table}

Two features of these results merit discussion. First, the test MAPE (17.42\%) is
lower than the training MAPE (28.56\%), an inverse of the usual pattern. This arises
because the training set includes the early 2018--2020 period when absolute EV sales
were very small---even modest absolute errors translate into large percentage errors
at low base levels. The test window (2023--2024) involves much higher volumes, so
the same absolute deviation represents a smaller percentage. Second, the test RMSE
(108,191) is roughly double the training RMSE (55,014), which is a modest and
acceptable train-to-test degradation for a highly volatile series. The most
challenging period in the test window is January--February 2024, when actual sales
fell sharply to around 290,000 units while the model predicted approximately 590,000,
reflecting the joint effect of Lunar New Year and post-year-end subsidy adjustments
that no univariate model could anticipate.

Figure~\ref{fig:test_forecast} plots the actual versus forecast values on the test set.

\begin{figure}[H]
\centering
\includegraphics[width=0.92\textwidth]{figures/08_test_forecast_nested_sarima.png}
\caption{Test set forecast: SARIMA$(1,1,0)\times(1,1,0)_{12}$ (red) versus actual EV
sales (blue), with 95\% confidence interval (shaded). The model tracks the seasonal
pattern well from May 2023 to December 2023 but underestimates the sharp early-2024
dip, which is attributable to the combined Lunar New Year and policy adjustment effect.}
\label{fig:test_forecast}
\end{figure}

\subsection{Robustness Check: Additive vs.\ Multiplicative Specification}
\label{sec:robust}

To test whether log-transforming the series (a common approach for variance
stabilization) improves forecasting, we re-run the full nested-CV pipeline on
$y^* = \log(1 + y_t)$ and back-transform the forecasts. Table~\ref{tab:log} reports
the comparison.

\begin{table}[H]
\centering
\caption{Additive (no-log) vs.\ multiplicative (log-space) SARIMA: test set accuracy}
\label{tab:log}
\begin{tabular}{lcc}
\toprule
Specification & Test RMSE (units) & Test MAPE (\%) \\
\midrule
Additive --- no-log (Nested CV) & 108,191 & 17.42 \\
Additive --- no-log (Auto-ARIMA) & 139,629 & 25.20 \\
Multiplicative --- log-space (Auto-ARIMA) & 322,831 & 56.02 \\
Multiplicative --- log-space (Nested CV) & 219,073 & 51.66 \\
\bottomrule
\end{tabular}
\end{table}

The log-space specifications perform dramatically worse in out-of-sample terms
(MAPE exceeding 50\%), despite appearing to fit the training data well. This
``explosive forecasting effect'' arises because the log model learns the rapid
growth phase of 2018--2022 and extrapolates it beyond the structural break of 2023.
When back-transformed, even modest log-space errors compound into very large
absolute unit errors. The result confirms that the additive SARIMA on the original
series is the correct model family for this dataset, particularly given the
post-subsidy structural change.

\subsection{Comparison with Baseline Models}

Table~\ref{tab:baselines} compares the selected SARIMA model against all baselines
on the same test set.

\begin{table}[H]
\centering
\caption{Forecast comparison on the test set (May 2023 -- April 2024)}
\label{tab:baselines}
\begin{tabular}{lcc}
\toprule
Model & Test RMSE (units) & Test MAPE (\%) \\
\midrule
SARIMA$(1,1,0)\times(1,1,0)_{12}$ (Nested CV) & 108,191 & \textbf{17.42} \\
Holt-Winters (ETS, additive) & \textbf{97,746} & 17.25 \\
Prophet & 119,070 & 18.45 \\
LSTM (lookback = 24) & 138,427 & 22.99 \\
ARIMA (non-seasonal) & 142,781 & 21.29 \\
\bottomrule
\end{tabular}
\end{table}

Holt-Winters achieves the lowest RMSE (97,746), slightly outperforming SARIMA on this
metric. This is because Holt-Winters happens to track the sharp January--February 2024
dip closely in this particular test window---the observations that dominate the squared
error metric. SARIMA is virtually tied on MAPE (17.42\% vs.\ 17.25\%), and unlike
Holt-Winters, its selection was validated across multiple time folds, not optimized
for a single test split. Prophet performs comparably but slightly worse than both
seasonal models (MAPE 18.45\%). The deep learning baseline (LSTM) underperforms all
classical methods, confirming that 76 monthly observations are insufficient for a data-hungry
architecture with 32 hidden units to learn the seasonal structure reliably. Non-seasonal
ARIMA is the worst performer (MAPE 21.29\%), confirming the importance of seasonal
components in this dataset.

\subsection{12-Month Forward Forecast}

The model is refitted on all 76 observations and used to forecast May 2024 to April 2025.
Table~\ref{tab:forecast} reports point forecasts and 95\% prediction intervals.

\begin{table}[H]
\centering
\caption{12-month forward forecast with 95\% prediction intervals}
\label{tab:forecast}
\begin{tabular}{lrrr}
\toprule
Month & Forecast (units) & Lower 95\% & Upper 95\% \\
\midrule
May 2024 & 536,581 & 397,903 & 675,259 \\
Jun 2024 & 637,615 & 473,204 & 802,026 \\
Jul 2024 & 612,811 & 429,239 & 796,382 \\
Aug 2024 & 677,390 & 479,336 & 875,444 \\
Sep 2024 & 697,665 & 486,770 & 908,559 \\
Oct 2024 & 699,527 & 476,915 & 922,140 \\
Nov 2024 & 753,361 & 519,740 & 986,981 \\
Dec 2024 & 804,413 & 560,341 & 1,048,485 \\
Jan 2025 & 522,333 & 268,262 & 776,405 \\
Feb 2025 & 446,472 & 191,170 & 718,534 \\
Mar 2025 & 624,180 & 351,229 & 897,130 \\
Apr 2025 & 576,284 & 294,371 & 858,196 \\
\bottomrule
\end{tabular}
\end{table}

\begin{figure}[H]
\centering
\includegraphics[width=0.92\textwidth]{figures/11_future_forecast.png}
\caption{12-month forward forecast (green line) with 95\% confidence interval
(shaded), plotted against the full historical series (blue). The forecast shows
a seasonal peak near December 2024 and a trough in February 2025. The confidence
interval widens as the horizon increases, reflecting cumulative forecast uncertainty.}
\label{fig:future_forecast}
\end{figure}

The forecast pattern is economically coherent and consistent with known Chinese EV
market dynamics. Sales rise steadily from May through December 2024, peaking at
approximately 804,000 units in December---driven by year-end promotional campaigns,
fleet purchase cycles, and the annual concentration of government-linked
institutional buying. A sharp seasonal correction is predicted in January--February
2025, with sales bottoming at roughly 446,000 units in February, as the Lunar New
Year holiday disrupts both production and retail. A recovery to the mid-600,000
range is forecast for March--April 2025 as market activity normalizes.

The 95\% confidence interval for December 2024 spans from 560,341 to 1,048,485
units, reflecting substantial but not unreasonable uncertainty at a 7-month horizon
in a rapidly evolving market. The widening interval over the forecast horizon is
consistent with theoretical expectations for ARIMA-family models \citep{chen2025,
cetin2024}.

%% ─────────────────────────────────────────────────────────────────────────────
\section{Discussion}
\label{sec:discussion}
%% ─────────────────────────────────────────────────────────────────────────────

\subsection{Economic Interpretation}

The SARIMA forecast captures two economically important features of China's EV
market that are consistent with the existing literature. First, the overall upward
trend---though muted relative to 2021--2022---reflects the continued baseline growth
in consumer EV adoption as vehicle prices decline, model variety increases, and
charging infrastructure expands. The China NEV penetration rate has risen from below
5\% to over 30\% of new passenger vehicle sales between 2020 and 2024, and
continuing structural demand supports a growth trajectory even in a post-subsidy
environment.

Second, the pronounced within-year seasonal cycle---with a Q4 peak and Q1 trough---is
a robust feature documented in both \citet{chen2025} and in the raw historical data.
The December peak specifically reflects a cluster of market-specific drivers: (a)
year-end purchase incentives offered by dealers and local governments, (b) completion
of annual fleet procurement by corporate and institutional buyers, and (c) consumer
behavior driven by annual bonus income concentration in Q4. The February trough reflects
the near-complete suspension of automotive retail activity during the Lunar New Year
holiday period, which typically encompasses 10--15 calendar days but affects consumer
activity for a broader 4--6 week window \citep{chen2025}.

These structural drivers are stable enough that a univariate seasonal model can
represent them reliably under normal market conditions. However, the model's failure
in January--February 2024---where actual sales collapsed further than the model
anticipated---illustrates the key limitation of univariate time series: it cannot
incorporate information about changes in government policy that shift the demand
baseline in a non-seasonal way.

\subsection{Practical Implications}

The forecast results have several actionable implications for industry participants
and policymakers.

\textbf{Production and inventory planning:} The December 2024 peak forecast of
approximately 804,000 units implies that manufacturers and their suppliers should
begin building finished-vehicle and battery inventory in September--October 2024 at
the latest to meet fourth-quarter demand. Conversely, the February 2025 trough
forecast provides a signal to reduce production rates and inventory builds in
January, avoiding stranded working capital during the low-demand Lunar New Year
period. The 95\% confidence intervals support scenario-based planning: a ``high
scenario'' based on the upper bound (1,048,485 units in December) can inform
maximum capacity utilization decisions, while the lower bound (560,341 units) can
serve as a conservative base case for financial planning.

\textbf{Charging infrastructure deployment:} The seasonally-adjusted forecast directly
informs the timing of public charging network expansion. Months forecasted above
700,000 units (August--December 2024) will generate substantially higher daily
charging demand, particularly in urban and highway corridor settings. Infrastructure
planners can use these projections to prioritize grid upgrade investments ahead of
high-demand periods, as highlighted by \citet{chen2025} in the context of China's
NEV infrastructure policy.

\textbf{Supply chain coordination:} Battery cell procurement---typically negotiated
3--6 months in advance---requires accurate rolling forecasts aligned with the
seasonal demand cycle. The model's forecast table enables procurement planners to
align lithium carbonate and cathode material orders with the mid-year (Jun--Aug)
and year-end (Nov--Dec) demand surges, reducing both shortage risk and excess
inventory cost.

\subsection{Why SARIMA is the Preferred Model}

Despite Holt-Winters achieving marginally lower RMSE on this particular test window,
SARIMA is the preferred model for three reasons. First, SARIMA's selection via nested
rolling CV provides a more robust validation than a single-split evaluation, making
it less sensitive to the specific characteristics of one 12-month hold-out period.
Second, SARIMA's parameters (AR, I, MA terms with seasonal analogs) have direct
economic interpretability: the negative non-seasonal AR coefficient reflects
short-term mean reversion in monthly sales, and the negative seasonal AR coefficient
reflects the year-over-year market stabilization since 2022. Holt-Winters' smoothing
parameters lack this structural interpretation. Third, SARIMA's explicit probabilistic
framework naturally provides confidence intervals that quantify forecast uncertainty
in a theoretically principled way, supporting the scenario-based planning described
above.

\subsection{Limitations}

This paper's results should be interpreted with three important caveats.

\textbf{Univariate model limitation:} SARIMA relies entirely on lagged values of the
target series. It cannot anticipate policy discontinuities (subsidy changes, purchase
restrictions), macroeconomic shocks (income growth, fuel price changes, exchange rate
movements), or competitive dynamics (new model launches, price wars). The January 2023
post-subsidy collapse and the early 2024 underperformance both demonstrate this
fundamental limitation.

\textbf{Short sample and structural breaks:} With only 76 monthly observations spanning
two qualitatively different market regimes (pre- and post-subsidy phase-out), the model
must simultaneously learn both the high-growth pattern of 2018--2022 and the more stable
pattern of 2023--2024. This dual-regime estimation may reduce the model's ability to
capture either regime cleanly, and SARIMA's assumption of stable seasonality over time
may not hold across the full sample period.

\textbf{Aggregation:} The analysis focuses on total EV sales without disaggregation
by vehicle segment (BEV vs.\ PHEV), price tier, or brand. BEVs and PHEVs have
distinct seasonal patterns and different sensitivities to policy changes; aggregate
forecasting blends these dynamics and may mask important sub-market trends.

%% ─────────────────────────────────────────────────────────────────────────────
\section{Conclusion}
\label{sec:conclusion}
%% ─────────────────────────────────────────────────────────────────────────────

This paper applies a fully documented SARIMA framework to forecast monthly EV sales
in China from a 76-month historical series spanning January 2018 to April 2024.
The primary methodological contribution is a two-stage model selection procedure
that combines BIC-guided auto-ARIMA with nested rolling cross-validation, producing
a final specification of SARIMA$(1,1,0)\times(1,1,0)_{12}$ that is more parsimonious
and more robust than the BIC-optimal in-sample model.

The key empirical findings are: (1) the China EV sales series is strongly
non-stationary in levels but achieves confirmed stationarity after one non-seasonal
and one seasonal difference, as validated by both ADF and KPSS tests; (2) the
selected SARIMA model achieves test-set MAPE of 17.42\% and passes the Ljung-Box
residual independence test at lag 12 ($p = 0.336$); (3) an additive specification
on the original series significantly outperforms a multiplicative log-space
alternative, due to the structural market shift in 2023; (4) the 12-month forward
forecast predicts a December 2024 peak of approximately 804,000 units and a
February 2025 trough of approximately 446,000 units, consistent with established
seasonal patterns in China's NEV market.

These forecasts have practical relevance for manufacturers, distributors,
policymakers, and infrastructure planners. The December peak projection supports
proactive Q3--Q4 inventory and production builds, while the February trough informs
planned production slowdowns and seasonal cash flow management. The 95\% confidence
intervals facilitate low/high scenario planning across stakeholder groups.

Future work should extend this framework in three directions. First, incorporating
exogenous policy variables---such as subsidy amounts, purchase restriction policies,
and credit availability---via SARIMAX would reduce the model's vulnerability to
policy shocks. Second, disaggregating the forecast by BEV versus PHEV segment, or
by major OEM, would provide more actionable granularity for industry decision-makers.
Third, the model should be updated on a rolling basis as new monthly observations
become available, with automated anomaly detection to flag structural breaks early.

%% ─────────────────────────────────────────────────────────────────────────────
%% APPENDICES
%% ─────────────────────────────────────────────────────────────────────────────
\appendix

\section{Implementation Details}
\label{app:code}

The entire pipeline was implemented in Python 3.11 using the following libraries:
\texttt{pandas} for data loading and monthly aggregation; \texttt{numpy} for
numerical operations; \texttt{matplotlib} and \texttt{seaborn} for visualization;
\texttt{statsmodels} for SARIMAX estimation, ADF, KPSS, and Ljung-Box testing;
\texttt{pmdarima} for auto-ARIMA and rolling cross-validation;
\texttt{prophet} for the Facebook Prophet baseline; and \texttt{tensorflow/keras}
for the LSTM baseline. The code is organized into two notebooks:
\texttt{01\_Data\_Prep\_and\_EDA.ipynb} (preprocessing, decomposition, outlier
analysis) and \texttt{02\_Time\_Series\_Modeling.ipynb} (stationarity tests, model
selection, diagnostics, forecast comparison). All random seeds are fixed at 42
for reproducibility.

\section{SARIMA Algorithm Outline}
\label{app:algorithm}

\noindent\textbf{Input:} Monthly EV sales series $\{y_t\}_{t=1}^{T}$, seasonal period $m=12$.

\noindent\textbf{Step 1 --- Data Preprocessing:} Remove duplicates; convert
\texttt{units\_sold} to integer; filter EV records; aggregate to monthly start
frequency; fill missing months with zero.

\noindent\textbf{Step 2 --- Exploratory Analysis:} Plot raw series, compute summary
statistics, perform additive decomposition to inspect trend, seasonal, and residual
components. Apply 3-sigma outlier check on residuals.

\noindent\textbf{Step 3 --- Stationarity:} Run ADF and KPSS on raw series and
differenced variants. Select $(d,D)=(1,1)$ based on dual test confirmation.

\noindent\textbf{Step 4 --- Initial Model Search (Stage 1):} Run auto-ARIMA with
$p,q \le 2$, $P,Q \le 1$, $d=1$, $D=1$, $s=12$, BIC criterion, stepwise search.
Record top-3 candidates by BIC.

\noindent\textbf{Step 5 --- Nested Rolling CV (Stage 2):} For each fold in a
rolling expanding window, re-fit auto-ARIMA and record one-step-ahead RMSE. Aggregate
mean RMSE across folds per order configuration. Select the order with lowest mean RMSE.

\noindent\textbf{Step 6 --- Residual Diagnostics:} Inspect residual ACF, histogram,
Q-Q plot, and Ljung-Box $Q(12)$ statistic.

\noindent\textbf{Step 7 --- Baseline Comparison:} Fit ARIMA, Holt-Winters, Prophet,
and LSTM on the same training split. Evaluate all models on the 12-month test set.

\noindent\textbf{Step 8 --- Final Forecast:} Refit selected SARIMA on all $T$
observations. Generate $h=12$ step-ahead forecasts with 95\% confidence intervals.

\noindent\textbf{Output:} SARIMA parameter estimates, test-set RMSE/MAPE, residual
test $p$-values, and the 12-month forecast table.

\section{Interpretation of SARIMA Parameters}
\label{app:parameters}

For the selected model SARIMA$(1,1,0)\times(1,1,0)_{12}$:
\begin{itemize}
\item $p=1$: One non-seasonal autoregressive term, capturing the mean-reverting
dependency of current monthly sales on the previous month's sales.
\item $d=1$: One non-seasonal differencing, removing the stochastic trend.
\item $q=0$: No non-seasonal moving average term is needed once AR and differencing
are included.
\item $P=1$: One seasonal autoregressive term at lag 12, capturing the year-over-year
dependency between same-month sales in consecutive years.
\item $D=1$: One seasonal differencing, removing the stochastic seasonal component.
\item $Q=0$: No seasonal moving average term is needed.
\item $s=12$: Annual seasonal period (monthly data).
\end{itemize}

The model is more parsimonious than the BIC-optimal auto-ARIMA selection
$(2,1,2)\times(1,1,1)_{12}$, with only 2 estimated parameters rather than 6, which
is appropriate for a training set of 64 observations. The negative sign of both AR
coefficients ($\phi_1 = -0.322$ and $\Phi_1 = -0.564$) implies mean reversion at
both the monthly and annual frequencies, consistent with a market that fluctuates
around a growing trend rather than exhibiting momentum-driven boom-bust cycles.

\section{Full Forecast Table}
\label{app:forecast_table}

\begin{longtable}{lrrr}
\caption{12-month forward forecast (May 2024 to April 2025) with 95\%
prediction intervals}
\label{tab:full_forecast} \\
\toprule
Month & Forecast (units) & Lower 95\% (units) & Upper 95\% (units) \\
\midrule
\endfirsthead
\toprule
Month & Forecast (units) & Lower 95\% (units) & Upper 95\% (units) \\
\midrule
\endhead
\bottomrule
\endfoot
May 2024 & 536,581 & 397,903 & 675,259 \\
Jun 2024 & 637,615 & 473,204 & 802,026 \\
Jul 2024 & 612,811 & 429,239 & 796,382 \\
Aug 2024 & 677,390 & 479,336 & 875,444 \\
Sep 2024 & 697,665 & 486,770 & 908,559 \\
Oct 2024 & 699,527 & 476,915 & 922,140 \\
Nov 2024 & 753,361 & 519,740 & 986,981 \\
Dec 2024 & 804,413 & 560,341 & 1,048,485 \\
Jan 2025 & 522,333 & 268,262 & 776,405 \\
Feb 2025 & 446,472 & 191,170 & 718,534 \\
Mar 2025 & 624,180 & 351,229 & 897,130 \\
Apr 2025 & 576,284 & 294,371 & 858,196 \\
\end{longtable}

%% ─────────────────────────────────────────────────────────────────────────────
%% REFERENCES
%% ─────────────────────────────────────────────────────────────────────────────
\newpage
\section*{References}
\addcontentsline{toc}{section}{References}

\begin{description}
\setlength\itemsep{0.4em}

\item \textbf{Chen, J., Huang, Y., Jin, H. and Chen, X.} (2025) `Research on the
forecasting of sales volume of new energy vehicles based on SARIMA model',
\textit{Journal of Education, Humanities and Social Sciences}, 55, pp.~76--82.

\item \textbf{\c{C}etin, B. and Ta\c{s}demir, \c{C}.} (2024) `Forecasting electric
vehicle sales using optimized SARIMA model: A two-year predictive analysis',
\textit{Veri Bilimi Dergisi}, 7(2), pp.~41--51.

\item \textbf{Felixzhao} (2024) \textit{China Automobile Monthly Sales Data (2018--2024.4)}.
Available at: \url{https://www.kaggle.com/datasets/felixzhao/china-automobile-monthly-sales-data-2018-2024-4}
(Accessed: May 2026).

\item \textbf{Neugeboren, R. and Jacobson, M.} (2014) \textit{Writing Economics: A Guide
for Harvard Economics Concentrators}. Harvard Writing Project. Minor revisions
October 2014.

\end{description}

\end{document}